\documentclass[11pt]{article}
\usepackage{preamble}
\renewcommand{\answer}[1]{}

\begin{document}

\ladderheading{AP Statistics \textperiodcentered\ Unit 3 \textperiodcentered\ Topic 3.13 \textperiodcentered\ B: 2026-12-17 \textperiodcentered\ E: 2027-01-29}{One Result, Two Claims}

\focusbanner{TODAY'S PROBLEM}

Suppose a utility has 500,000 online accounts. It randomly selects 10,000, randomly assigns 5,000 to message A and 5,000 to B, and records on-time payment. The pre-data question asks whether the population proportions differ at $\alpha=0.05$; deployment requires A to improve payment by at least 5 percentage points.\par\smallskip \textbf{Rowan:} ``I rejected equality, so A has proved an important improvement of at least 5 points. Deploy A.''\par \textbf{Salma:} ``Rejecting equality supports a nonzero difference, but does not automatically establish the separate 5-point requirement.''\par\smallskip Conditions have been verified: random selection and assignment; $10{,}000\leq50{,}000$; pooled expected counts 3,000 and 2,000 in each group.\par
\begin{center}
\textbf{On-time payment}\par\smallskip
\begin{tabular}{|l|c|c|c|}
\hline
\textbf{Text} & \textbf{n} & \textbf{Paid} & \textbf{Other} \\ \hline
\textbf{A} & 5000 & 3050 & 1950 \\ \hline
\textbf{B} & 5000 & 2950 & 2050 \\ \hline
\end{tabular}
\end{center}
\begin{firsttakebox}{FIRST TAKE --- before the video (5 min)}
The observed improvement is 2 percentage points and the deployment target is 5. Is that enough to deploy? Give one reason.
\workspace{0.9in}
\end{firsttakebox}

\begin{callout}{\IconBook\ CALCULATE, INTERPRET, THEN BOUND THE CLAIM (keep on the page)}{calloutgreen}
For $H_0:p_1-p_2=0$, pool $\hat p_c=(x_1+x_2)/(n_1+n_2)$ and use
$z=\dfrac{(\hat p_1-\hat p_2)-0}{\sqrt{\hat p_c(1-\hat p_c)(1/n_1+1/n_2)}}$.
A two-sided test uses $P(|Z|\ge |z|)$. The $p$-value assumes $H_0$ is true and measures
how often random chance would produce a difference at least as extreme. If the $p$-value is at most $\alpha$, reject
$H_0$; otherwise fail to reject it. A test against zero addresses whether a difference exists;
it does not by itself establish that the difference exceeds a separate practical threshold.
\end{callout}

\newpage
\textbf{Finish after the video:}\par\smallskip

\dokbadge{1}~\textbf{(a)} Calculate $\hat p_A$, $\hat p_B$, and the observed difference $\hat p_A-\hat p_B$.\par
\workspace{0.7in}

\dokbadge{2}~\textbf{(b)} Using the verified conditions, calculate the pooled proportion, pooled SE, $z$, and the two-sided $p$-value for $H_0:p_A-p_B=0$.\par
\workspace{1.4in}

\dokbadge{3}~\textbf{(c)} In 3--4 sentences, choose a report using the $p$-value, the 2-point observed gap, and the 5-point target. Explain what the test against zero establishes and why that does or does not justify deployment.\par
\workspace{2.0in}

\begin{sentenceframebox}\raggedright\textbf{Frame:} The p-value supports \blank[2.0in], because \blank[1.8in].\end{sentenceframebox}

\begin{sentenceframebox}\raggedright\textbf{Frame:} The test compares the difference with \blank[0.7in], but deployment requires \blank[1.1in]; therefore \blank[2.0in].\end{sentenceframebox}

\begin{summaryexitbox}{EXIT --- turn this in}
One thing the video changed about my first take: \hrulefill\par
\workspace{0.4in}
\end{summaryexitbox}

\end{document}
